\section{Current Evidence and Evaluation Plan}
\label{sec:evaluation}
\subsection{Implemented artifact and local observations}
The implementation baseline is commit \implementationcommit{} on the QUIC
branch. It adds the packet/audio adaptation while reusing the existing modem,
identity, gateway, AVA, PAP, risk, and provider components. The complete local
Python regression suite reports \testcount{} passing tests. This count is not
a claim of security coverage completeness or a substitute for measurements.
The experiment inputs and sanitized reports are in
\path{docs/quic-verification.json}; Table~\ref{tab:local} is generated directly
from that file.

\begin{table}[t]
\centering
\begin{tabular}{llrr}
\toprule
Local audio route & Faults & Handshake (s) & $n$ \\
\midrule
\measurementrows
\bottomrule
\end{tabular}
\caption{Observed caller authentication time, including the QUIC server's mutual-authentication confirmation; excludes enrollment and call setup. Each row is one local feasibility run, not a latency distribution. All rows complete one allowed and seven denied probes.}
\label{tab:local}
\end{table}

All runs use separate caller and business Python processes on one macOS host
(arm64, Python 3.13.7, aioquic 1.3.0, aiortc 1.15.0), with locally provisioned
identities. WebRTC carries G.711 PCMU audio only. The provider-emulation route
uses the Twilio adapter and local media/playback envelopes, without a carrier
call. The lossy PCM runs drop two modem bursts, damage one, and duplicate one
in each direction using fixed burst indices. Native QUIC records three
packet losses at the caller and four at the business in those runs. Modem CRC
failures confirm that corruption reached the receiver. These are specific
waveform faults, not a claimed 1\% or 5\% network loss rate.

Each SDK call allows an order read, then rejects scope escalation, replay,
another customer, an expired grant, a revoked grant, another key, and a revoked
certificate. Each business reports exactly one handler invocation and seven
denials. The FSK run has no injected loss and completes the same probes; its
single observed handshake is approximately 35.7 seconds. A separate lossy
PAP/WebRTC run verifies a permitted read, an exchange proposal followed by
one HTTPS-confirmed simulated effect, repeated confirmation/status recovery,
fraud blocking, and denial after sign-out. It is a local reference service,
not interoperability with an external PAP provider.

Regression tests also cover streaming FSK samples, fragment reordering and
duplicates, reassembly bounds and expiry, corrupted ciphertext with a valid
audio CRC, large messages, malformed application input, ambiguous execution
outcomes, negative certificates, and final ACK playback. We have not measured
commercial detector accuracy, real-user approval, real telephone routes, or
physical microphone operation. Historical repository TLS timings on another
host and date are not treated as a controlled TLS-versus-QUIC comparison.

\subsection{Research questions and baselines}
The remaining evaluation is a prospective plan. We define six questions with
outcomes that can falsify a deployment claim:
\begin{enumerate}
\item \textbf{Authorization:} Does the gateway prevent unauthorized effects when
models encounter malicious speech, valid but insufficient credentials, replay,
handoff, or permission changes? Count actual backend effects as well as response codes.
\item \textbf{Transport:} Under which rates, losses, and delays does native QUIC
improve completion time or reliability over the existing reliable-audio TLS profile?
\item \textbf{Path compatibility:} Which real codec chains and processing settings
support each modem, and where does the system reliably fail closed?
\item \textbf{Deployment value:} When does in-audio control offer an advantage
over an authenticated API or an HTTPS grant correlated with a call?
\item \textbf{Conversation cost:} How much speech delay, interruption, audibility,
and user confusion do modem bursts introduce?
\item \textbf{Resource cost:} How do adversarial pre-authentication traffic,
many calls, and stale partial packets affect CPU, memory, and availability?
\end{enumerate}

Transport baselines must use the same modem, certificates, payloads, call
adapters, and host configuration: (a) current stop-and-wait TLS, (b) native
QUIC, and (c) an improved reliable audio link with a window or selective repeat.
The third baseline tests whether a benefit comes merely from replacing
stop-and-wait. An authenticated HTTPS API supplies a deployment baseline;
an HTTPS grant plus a securely specified call-binding mechanism supplies a
cross-channel baseline. A raw unauthenticated speech interface is a negative
control for authorization, not a competitive secure protocol.

\subsection{Channel and performance experiments}
Start with lossless sample loopback, local WebRTC, and provider emulation to
locate errors before introducing uncontrolled routes. Then use consented
endpoints across controlled Internet links, real provider calls, and finally
microphone/speaker paths in a laboratory. Log codec negotiation and actual
resampling, echo cancellation, automatic gain control, noise suppression,
packetization, and jitter-buffer behavior. Test G.711, representative
transcoding chains, Opus settings, clock mismatch, sample insertion/deletion,
gain, noise, clipping, and speech overlap. A codec/profile combination that
cannot decode should be reported as unsupported rather than excluded from
the results silently.

For digital impairments, sweep independently seeded packet loss at 0, 0.1, 1,
3, 5, and 10\%, and include a burst-loss model with stated parameters. Sweep
one-way network delay, jitter, reordering, and duplicate packets separately
from waveform corruption and modem-fragment erasure. These fault units are
not interchangeable. Use at least 100 fresh calls per selected condition
initially, then increase sample sizes where tail estimates are unstable.
With zero failures in 100 trials, a rough 95\% upper bound on failure
probability is still about 3\%; do not interpret such a result as five-nines
reliability. Report binomial intervals, medians and tail confidence intervals,
timeouts, and all failed calls.

Measure complete mutual-authentication time at \emph{both} endpoints, first
authorized operation, transaction completion, airtime, useful goodput,
retransmitted bytes, CPU, peak memory, and speech-queue delay. Separate call
setup, enrollment, grant acquisition, handshake, application exchange, and
LLM/speech times. Randomize baseline order and pair runs by impairment seed.
Keep raw timings and wall-clock/software metadata. Sweep certificate-chain
size, request/response size, audio frame size, initial RTT, and burst limits.
Do not introduce resumption or 0-RTT into one baseline alone.

\subsection{Adversarial and compositional experiments}
Build an attack matrix for wrong roots, service names and purposes, missing
certificates, invalid key proof, stolen grants under another certificate,
concurrent grant use, stale nonces, argument substitution, revocation, and
disconnects at every execution boundary. Exercise duplicate keys, multiple
calls sharing an identity, partial authentication, attempted downgrade, live
relay, and transfers between endpoints. Relays are expected to defeat a
physical-presence interpretation; a correct evaluation should expose that
limit rather than count successful authentication as relay prevention.

For prompt injection, compare identical model tasks with and without gateway
enforcement. Include in-scope malicious proposals, out-of-scope proposals,
speech-only attempts, and responses that try to extract secrets or trigger
another tool. Record model proposals, gateway decisions, real effects, and
leaks through all enabled outputs. Fix and report the model version and prompts,
and use repeated independent trials. A blocked financial write does not
establish that the agent retained unrelated confidential information.

Test failures during an asynchronous risk check and between revalidation and
commit. Inject grant-service outages, risk-provider errors, and provider
disconnects, checking bounded resource use and absence of silent fallback.
Audit keys and tokens in logs, traces, model inputs, ASR output, and recordings.
Fuzz fragment headers and authenticated application messages. A symbolic model
of the grant/challenge state machine and an independent review of private
aioquic hooks are planned work; neither is claimed complete here.

\subsection{Users, operations, and decision criteria}
Any user study should use informed consent and a clear transaction-approval
interface. Measure whether participants understand the authenticated principal,
the approved action, speech outside protection, and what happens on a transfer
or timeout. Measure business integration effort and operator handling of
unsupported calls. Publish the intervention and limitations, not just task
completion. An operational recommendation requires both acceptable authorized
completion on stated paths and reliable rejection under the attack matrix.
If a conventional authenticated API achieves the required binding with less
integration and user cost, the evaluation should recommend it.
